\subsection{Existence of simple minimizers}
\label{subsec:generalized-reeb-orbits-polytopes-existence-of-simple-minimizers}

Haim--Kislev prove that every convex polytope has a minimum-action generalized
Reeb orbit with a simple finite description \cite[Theorem~1.5]{HK2017}. In
their statement the velocities are positive multiples of \(J_0n_i\) for the
unit facet normals \(n_i\). Here we reparametrize each segment by its
contact-normalized facet direction \(R_i=2J_0a_i\), absorbing the positive
multiples into the dwell times. The proof is important for the finite
computation, so we record the argument in the conventions of this thesis. It
starts from an arbitrary minimizer, passes to the dual formulation from
Theorem~\ref{thm:preliminaries-clarke-dual-action}, and constructs simpler
dual curves whose values converge back to the minimum.

\begin{theorem}[Simple minimizer]\label{thm:generalized-reeb-simple-minimizer}
Let \(K\subset\mathbb{R}^4\) be a full-dimensional convex polytope with
\(0\in\operatorname{int}(K)\). Among all generalized Reeb orbits on
\(\partial K\), there exists a minimum-action orbit that is simple in the
contact-normalized convention above. Hence
\[
  c_{\mathrm{EHZ}}(K)
  =
  \min\{A(\gamma):\gamma
  \text{ is a simple Reeb orbit on }\partial K\}.
\]
\end{theorem}

The proof uses five elementary operations on dual curves. The first operation
replaces a \(W^{1,2}\)-curve by a piecewise affine one.

\begin{lemma}[Piecewise affine approximation]\label{lem:generalized-reeb-pw-affine-approx}
Let \(z\in W^{1,2}([0,T],\mathbb{R}^4)\) be closed and satisfy
\[
  \dot z(t)\in\operatorname{conv}\{R_1,\ldots,R_F\}
  \quad\text{for a.e. }t.
\]
For every \(\varepsilon>0\), there is a closed piecewise affine curve \(z'\)
with
\[
  \dot z'(t)\in\operatorname{conv}\{R_1,\ldots,R_F\}
\]
on each affine piece and \(\|z'-z\|_{W^{1,2}}<\varepsilon\). Consequently
\(A(z')\to A(z)\) and \(I_K(z')\to I_K(z)\) as \(\varepsilon\to0\).
\end{lemma}

\begin{proof}
Let \(C=\operatorname{conv}\{R_1,\ldots,R_F\}\). Conditional expectations onto
refining finite partitions converge to \(\dot z\) in \(L^2\). Choose one such
partition \(0=t_0<\cdots<t_m=T\) so that the conditional averages of
\(\dot z\) on the partition intervals are \(L^2\)-close to \(\dot z\). On
\([t_k,t_{k+1}]\), put
\[
  w_k=\frac{1}{t_{k+1}-t_k}\int_{t_k}^{t_{k+1}}\dot z(s)\,ds.
\]
Because \(C\) is convex and closed, each \(w_k\) lies in \(C\). Define \(z'\)
by \(z'(0)=z(0)\) and \(\dot z'=w_k\) on \([t_k,t_{k+1}]\). Then
\[
  z'(T)-z'(0)=\int_0^T\dot z'(t)\,dt=\int_0^T\dot z(t)\,dt=0,
\]
so \(z'\) is closed. Since \(\dot z'\to\dot z\) in \(L^2\), integration gives
\(z'\to z\) uniformly, hence in \(W^{1,2}\). The convergence of \(A\) follows
from bilinearity in \((z,\dot z)\), and the convergence of \(I_K\) follows
from continuity of \(h_K^2\) on the bounded velocity set.
\end{proof}

\begin{lemma}[Splitting mixed velocities]\label{lem:generalized-reeb-splitting}
Let \(z\) be a closed piecewise affine curve of period \(T\) whose velocities
lie in \(\operatorname{conv}\{R_1,\ldots,R_F\}\). Then there is a closed
piecewise affine curve \(z'\), with the same period, whose velocity on each
segment is one pure Reeb vector \(R_i\), such that \(A(z')\geq A(z)\) and
\[
  I_K(z')=T.
\]
\end{lemma}

\begin{proof}
On a segment with velocity \(v=\sum_i\theta_iR_i\), \(\theta_i\ge0\) and
\(\sum_i\theta_i=1\), replace the segment by subsegments with velocities
\(R_i\) and lengths \(\theta_i\) times the original length. The total
displacement and period are unchanged. The only change in the action is the
sum of the new internal cross-terms from
Lemma~\ref{lem:preliminaries-shoelace}. Reversing the order of the inserted
subsegments changes the sign of this internal contribution, so one of the two
orders does not decrease the action.

For a pure Reeb velocity \(R_i=2J_0a_i\),
\[
  -J_0R_i=2a_i,
  \qquad
  h_K(-J_0R_i)=2h_K(a_i)=2,
\]
because \(F_i\neq\varnothing\) and \(\max_{x\in K}\langle a_i,x\rangle=1\).
Thus \(h_K^2(-J_0R_i)=4\) on every pure segment, and
\[
  I_K(z')=\frac14\int_0^T4\,dt=T.
\]
\end{proof}

\begin{lemma}[Merging repeated velocities]\label{lem:generalized-reeb-merging}
Let \(z\) be a closed piecewise affine curve whose velocity on each segment is
one of the \(R_i\). If some \(R_i\) occurs in two non-adjacent segments, then
there is a rearrangement \(z'\) that merges those two occurrences into one
contiguous segment and satisfies
\[
  A(z')\geq A(z),
  \qquad
  I_K(z')=I_K(z).
\]
Repeating this operation, and merging adjacent equal-velocity intervals,
gives a closed piecewise affine curve in which each pure Reeb vector occurs in
at most one interval. Here adjacency is understood cyclically, since the curve
is closed.
\end{lemma}

\begin{proof}
Suppose a velocity \(v\) occurs in two segments \(I_r\) and \(I_s\), with
middle segments \(M=\{r+1,\ldots,s-1\}\). Merging \(I_s\) forward into \(I_r\)
changes only the action terms involving \(v\) and the middle velocities, and
the change is
\[
  \Delta A_{\mathrm{fwd}}
  =
  |I_s|\sum_{i\in M}|I_i|\,\omega_0(v,w_i).
\]
Merging \(I_r\) backward into \(I_s\) gives
\[
  \Delta A_{\mathrm{bwd}}
  =
  -|I_r|\sum_{i\in M}|I_i|\,\omega_0(v,w_i).
\]
The two changes have opposite signs, or both vanish, so one merge has
nonnegative action change. Closure is preserved because the multiset of
velocity-duration pairs is unchanged, and \(I_K\) is unchanged because it
depends only on the velocities and their durations.
\end{proof}

\begin{lemma}[Rescaling]\label{lem:generalized-reeb-rescaling}
Let \(z\) be a closed piecewise affine curve with pure Reeb velocities and
period \(T\). If all durations are multiplied by \(\beta>0\), the resulting
curve \(z'\) satisfies
\[
  A(z')=\beta^2A(z),
  \qquad
  I_K(z')=\beta I_K(z).
\]
\end{lemma}

\begin{proof}
In Lemma~\ref{lem:preliminaries-shoelace}, every action term contains a
product of two durations, so \(A\) scales by \(\beta^2\). The integrand in
\(I_K\) is unchanged and the time intervals are all multiplied by \(\beta\), so
\(I_K\) scales by \(\beta\).
\end{proof}

\begin{lemma}[Compactness of simple dual curves]\label{lem:generalized-reeb-compactness}
Let \(z_n\) be closed simple piecewise-Reeb curves whose periods satisfy
\(T_n\to T>0\). After translating the curves and passing to a subsequence,
\(z_n\) converges to a closed simple piecewise-Reeb curve \(z_*\) with period
\(T\), and
\[
  A(z_*)=\lim_n A(z_n),
  \qquad
  I_K(z_*)=T.
\]
\end{lemma}

\begin{proof}
Choose the initial point of each periodic curve at a breakpoint, and insert
zero dwell times for unused facet directions. Each curve is then determined up
to translation by a permutation of the finitely many facet directions, dwell
times in the simplex of total mass \(T_n\), and its period. Since
the periods eventually lie in a compact interval \([T/2,2T]\), a subsequence of
these finite-dimensional data converges. Translating the starting points gives
uniform convergence after identifying each domain with \([0,1]\) by linear
time rescaling. Closure passes to the limit because
\(\sum_i\tau_iR_{\sigma(i)}=0\) for every curve. Some limiting dwell times may be zero; deleting
those zero-duration entries recovers the active word of the limiting simple
curve. The action is continuous in the dwell-time data by
Lemma~\ref{lem:preliminaries-shoelace}. Since the remaining positive-duration
segments still have pure Reeb velocities, the same computation as in
Lemma~\ref{lem:generalized-reeb-splitting} gives \(I_K(z_*)=T\).
\end{proof}

\begin{proof}[Proof of Theorem~\ref{thm:generalized-reeb-simple-minimizer}]
Choose a capacity-realizing generalized Reeb orbit \(\gamma\). By
Theorem~\ref{thm:preliminaries-clarke-dual-action}, \(\gamma\) is also a dual
minimizer in the thesis convention, with
\[
  I_K(\gamma)=A(\gamma)=T=c_{\mathrm{EHZ}}(K).
\]
Moreover \(\dot\gamma(t)\in\operatorname{conv}\{R_1,\ldots,R_F\}\) for almost
every \(t\).

Apply Lemma~\ref{lem:generalized-reeb-pw-affine-approx} to approximate
\(\gamma\) by closed piecewise affine dual curves \(z_1\) with velocities in
the convex hull of the \(R_i\). Then apply
Lemma~\ref{lem:generalized-reeb-splitting} to obtain a closed pure-velocity
curve \(z_2\) with \(A(z_2)\geq A(z_1)\) and \(I_K(z_2)=T\). Apply
Lemma~\ref{lem:generalized-reeb-merging} to obtain a closed pure-velocity curve
\(z_3\) in which each \(R_i\) occurs at most once, still with
\[
  A(z_3)\geq A(z_2),
  \qquad
  I_K(z_3)=T.
\]
We take the initial approximation close enough that \(A(z_1)>T/2\); then
\(A(z_3)>0\). Now rescale all durations by
\[
  \beta=\frac{T}{A(z_3)}.
\]
The rescaled curve \(z_4\) satisfies
\[
  A(z_4)=\frac{T^2}{A(z_3)},
  \qquad
  I_K(z_4)=\frac{T^2}{A(z_3)}.
\]
Thus \(z_4\) is feasible for the dual problem with period \(T^2/A(z_3)\).
Since \(\gamma\) was a dual minimizer, \(I_K(z_4)\geq T\), and therefore
\(A(z_3)\leq T\). On the other hand \(A(z_3)\geq A(z_1)\), and
\(A(z_1)\to T\). Hence \(A(z_3)\to T\), the rescaling factors tend to \(1\),
and the periods of \(z_4\) tend to \(T\).

By Lemma~\ref{lem:generalized-reeb-compactness}, a subsequence of the \(z_4\)
converges up to translation to a closed simple piecewise-Reeb curve \(z_*\)
with \(A(z_*)=T\) and \(I_K(z_*)=T\). This curve is a dual minimizer. Applying
Theorem~\ref{thm:preliminaries-clarke-dual-action} again gives a
contact-normalized generalized Reeb orbit on \(\partial K\) in the same
translate/rescale family. In the reconstruction part of that theorem, the
corresponding characteristic has the form
\[
  t\longmapsto c\,z_*(t/c^2)+b/c
\]
for some \(c>0\), with
\[
  c^2=\nu=\frac{I_K(z_*)}{T}=1
\]
by \eqref{eq:preliminaries-dual-multiplier-value}. Thus the reconstruction
only translates \(z_*\). Its velocities
remain the pure facet directions \(R_i\), each used in one connected interval,
so the resulting orbit is simple and has minimum action.
\end{proof}

The theorem proves existence of at least one simple capacity minimizer. It is
not a classification of all minimum-action generalized Reeb orbits. Other
minimizers may pass through nonsmooth faces and may use convex combinations of
active facet directions. Thus the finite Haim--Kislev quadratic program is
justified as a capacity search because its search family contains a simple
capacity minimizer, not because every minimizer or every word has this form.
Infeasible words and duplicate descriptions can still occur.
